%---------------------------------------------------------------------
% Appendix: ECSA Graduate Attributes
%---------------------------------------------------------------------
\chapter{ECSA Graduate Attributes}
\label{app:ga}

This appendix shows how each of the six Engineering Council of South Africa (ECSA) Graduate Attributes (GAs) assessed in EEE4022S was met. Each GA is covered in its own section below. The plan for each GA was recorded in the GA tracking form submitted on 18 September 2026.

%---------------------------------------------------------------------
\section{GA1: Problem Solving}
\label{app:ga1}
%---------------------------------------------------------------------

The problem was not clearly defined at the start. The topic asked for an incubator that controls temperature, humidity and CO\textsubscript{2}, but the biology gave no exact tolerances, and the bacterium does not even need added CO\textsubscript{2} (Section~\ref{subsec:myco-co2-humidity}). The problem was solved in the following steps.

\begin{itemize}[nosep]
    \item \textbf{Defining the problem.} The problem was stated without assuming a solution (Section~\ref{sec:problem-statement}) and turned into ten requirements, R1 to R10, each with a stated reason (Section~\ref{sec:requirements}).
    \item \textbf{Balancing conflicts.} Some requirements worked against each other. The heated lid stops condensation but sits in the light path, and heating the plate from below warms the sample but adds a surface under the objective. Each trade-off was decided and recorded (Sections~\ref{sec:enclosure}, \ref{sec:temperature-design} and~\ref{sec:cl-test3}).
    \item \textbf{Using first principles.} Heat transfer, mass balance and vapour pressure physics were used to model the chamber. The chamber and glass were modelled from energy balances, and the gas from a mass balance (Chapter~\ref{ch:theory}). The controllers were designed from these models so that they would not overshoot, since the system can only heat.
    \item \textbf{Correcting from measurements.} The measured gains showed that the two temperature loops affect each other strongly, so the chamber controller was retuned with a corrected gain of $K_\mathrm{eff} = 32.6$, which was confirmed in testing (Section~\ref{sec:sysid-gains}).
    \item \textbf{Finding the causes of failures.} The cold sample was traced to how the chamber is heated (Chapter~\ref{ch:ClosedLoop}), and the humidity limit was calculated from the water and air temperatures, matching the measured 78\,\% (Section~\ref{sec:humidity}). Both point to one fix, heating the plate from below, which is the first item in Chapter~\ref{ch:furtherwork}.
    \item \textbf{Wider factors.} The build cost R\,13\,527, far below commercial systems costing several hundred thousand rand. A retail CO\textsubscript{2} cylinder was used so that others can rebuild it easily, and the CO\textsubscript{2} channel was kept for later BSL-3 work. This keeps the design affordable and possible to rebuild for groups with limited funding.
\end{itemize}

%---------------------------------------------------------------------
\section{GA4: Investigations, Experiments and Data Analysis}
\label{app:ga4}
%---------------------------------------------------------------------

\begin{itemize}[nosep]
    \item \textbf{Literature review.} Chapter~\ref{ch:review} covered the biology, nine existing incubators and the control of coupled loops. It led directly to design choices: keeping CO\textsubscript{2}, using separate PI loops, and using a cylinder and solenoid.
    \item \textbf{Design process.} The design went from requirements, to comparing options (Appendix~\ref{app:system-design}), to block diagrams (Figures~\ref{fig:block-diagram}, \ref{fig:temp-diagram} and~\ref{fig:co2-subsystem}), to testing each part on its own (Chapter~\ref{ch:commissioning}), and then testing the full system (Chapters~\ref{ch:thermalID} to~\ref{ch:co2humidity}).
    \item \textbf{Experiments.} The PT1000 was checked in two separately made ice baths, to tell an error in the ice apart from an error in the measurement, and the thermistors were calibrated against it (Section~\ref{subsec:pt1000-commissioning}). Step tests measured each temperature loop and the coupling between them (Chapter~\ref{ch:thermalID}). Four closed-loop runs tested holding the setpoint, warming the sample, heating from below and evenness across the plate (Chapter~\ref{ch:ClosedLoop}). A 6.5 hour run measured humidity (Section~\ref{sec:humidity}). When the gas joint failed, a sensitivity study was done instead and recommended what to measure first in the next build (Section~\ref{sec:co2-sensitivity}).
    \item \textbf{Data analysis.} Uncertainties were given for the measured time constants, and gains were taken in a way that cancels out room temperature changes. The 0.038\,\textdegree{}C standard deviation in closed loop was shown to be set by the sensor resolution, not the controller (Section~\ref{sec:cl-test1}). Every result was reported against the sensor limits, and every requirement was checked against the evidence, including those not met (Table~\ref{tab:disc-requirements}).
    \item \textbf{Code.} The firmware and analysis code are in the repository (Appendix~\ref{app:github}), and the firmware was tested one channel at a time before the system was put together (Chapter~\ref{ch:commissioning}).
\end{itemize}

%---------------------------------------------------------------------
\section{GA5: Use of Engineering Tools}
\label{app:ga5}
%---------------------------------------------------------------------

The tools used were:

\begin{itemize}[nosep]
    \item \textbf{MATLAB and Simulink} for modelling, controller design, robustness checks and frequency analysis (Chapter~\ref{ch:theory}, Appendix~\ref{app:control-results}), and for fitting models to the step test data (Chapter~\ref{ch:thermalID}).
    \item \textbf{Embedded C++ on the STM32} in PlatformIO, for reading the sensors over I\textsuperscript{2}C, SPI and UART, driving the heaters and solenoid, running the controllers and logging data (Sections~\ref{sec:power-design} and~\ref{sec:firmware-settings}).
    \item \textbf{CAD and 3D printing} in PETG for the enclosure, designed to measurements taken from the microscope stage (Section~\ref{sec:enclosure}).
    \item \textbf{Circuit design} for the heater drivers, thermistor dividers and the relay that switches the 230\,V solenoid (Figure~\ref{fig:driver-circuits}).
    \item \textbf{Lab methods}: multimeter checks, ice bath calibration, the CO\textsubscript{2} sensor's fresh-air zero, and a soapy-water leak test (Chapter~\ref{ch:commissioning}, Section~\ref{sec:co2c-failure}).
    \item \textbf{Standards}: IEC 60751 for RTD accuracy, ANSI/SLAS 1-2004 for the plate size, and ASTM D6245 and E741 for the gas model and leak test.
    \item \textbf{\LaTeX, BibTeX, Git and GitHub} for the report and version control.
\end{itemize}

The limits of these tools were stated where they mattered. The models treat each part as one temperature, which the cold corner of the enclosure shows is not always true (Section~\ref{sec:disc-models}). The controller tuning assumes the model is exact, so robustness was checked in simulation (Section~\ref{subsec:tuning}). The sensors limit how precisely temperature can be stated, so no claim was made beyond their accuracy (Section~\ref{sec:disc-requirements}).

%---------------------------------------------------------------------
\section{GA6: Professional and Technical Communication}
\label{app:ga6}
%---------------------------------------------------------------------

The report follows the department's template and the order in which the work was done, so results always come after the design they depend on. Equations are typed and numbered, figures have captions below and tables have captions above, and every figure and table is referred to in the text. All sources are cited in IEEE style, including figures taken from other work. Acronyms are defined when first used and listed in the front matter with the symbols. The report is written in the third person, and requirements that were not met are stated clearly with the reason (Table~\ref{tab:disc-requirements}). The system is shown with block diagrams, circuit diagrams and CAD drawings. The report is written for a mixed audience: the biological background is explained for engineers, and the main results are clear enough for a non-engineer to follow.

%---------------------------------------------------------------------
\section{GA8: Individual, Team and Multidisciplinary Working}
\label{app:ga8}
%---------------------------------------------------------------------

\begin{itemize}[nosep]
    \item \textbf{Individual work.} The project was done individually, with guidance from the supervisor, R.A. Verrinder, and co-supervisor, Dr M.K. Mason of the Molecular Mycobacteriology Research Unit.
    \item \textbf{Multidisciplinary work.} The requirements came from microbiology and were set through regular discussion with the co-supervisor, who works in microbiology. The design used thermal modelling, gas handling, electronics, control and embedded systems. Suppliers were also consulted directly, for example to confirm the regulator fitting before ordering (Section~\ref{sec:co2-design}).
    \item \textbf{Ethics.} The Faculty ethics pre-screening was approved on 11 August 2026 (reference \ebeethicsnumber), before any testing. The letter is in Appendix~\ref{app:ethics}.
    \item \textbf{Safety.} The mains solenoid was switched through a relay, with the mains wiring kept in a closed box (Section~\ref{sec:power-commissioning}). A washer that would have sealed the leaking gas joint was not used, because it left only one thread holding against 57\,bar (Section~\ref{sec:co2c-failure}).
    \item \textbf{Honesty.} Results were reported as measured, including the requirements that were not met, and all help, including the use of AI tools, is acknowledged in the Acknowledgements and the Plagiarism Declaration and described in Appendix~\ref{app:ai}.
\end{itemize}

%---------------------------------------------------------------------
\section{GA9: Independent Learning Ability}
\label{app:ga9}
%---------------------------------------------------------------------

Most of what this project needed is not taught in the undergraduate curriculum and was learnt independently:

\begin{itemize}[nosep]
    \item the biology of mycobacteria and what live-cell imaging needs;
    \item thermal modelling of the chamber and glass, and modelling of the gas in the chamber, neither of which is covered in the courses;
    \item compressed gas components, including the cylinder, regulator, solenoid valve, needle valve and fittings, none of which had been taught;
    \item sensors and components new to the project, including how NDIR CO\textsubscript{2} sensors work and are zeroed, PT1000 accuracy classes and lead-resistance error, and the ITO heated glass;
    \item measuring a system's response, and how coupling between loops changes the controller tuning;
    \item the physics that limits the humidity in the chamber.
\end{itemize}

Much of the project involved working without clear answers. The biology gave no exact tolerances, so the requirements had to be set and justified without a fixed target. Where knowledge was lacking, help was asked for: from Dr Mason on the biology, from suppliers on the gas fittings and parts, and from R.A. Verrinder in weekly meetings, whose feedback and advice on the design, testing and report were used to improve the work as it went. When the humidity result fell short, the physics behind it was learnt and used to explain it. The single-point calibration and the unchecked humidity sensors were recognised as gaps and are listed as further work.

The project was managed through these weekly meetings and a milestone plan. Several problems were solved independently along the way. A reversed power connection damaged the RTD converter board, and the fault was traced to its regulator and repaired (Section~\ref{subsec:pt1000-commissioning}). The relay would not switch off from 3.3\,V logic, which was fixed by setting the pin to open-drain (Section~\ref{sec:power-commissioning}). When the gas joint failed the day before testing, the cause was found and the CO\textsubscript{2} work was changed to a sensitivity study (Section~\ref{sec:co2c-failure}). The humidity work was likewise changed to a passive test that still answered a useful question. The report ends with an honest check against every requirement (Section~\ref{sec:disc-requirements}) and a list of what the next build should change (Chapter~\ref{ch:furtherwork}).
